\cotefile{156-5}
\scanpage{26}\mspage{25}
% @ top, lines 1–4
(\unsure{resp.} \unsure{ouvert}, \unsure{resp.} \unsure{loc. fer.}) est \unsure{idem}. Le cas d'une \deleted{\unsure{sur|su}}\unsure{subdivision} est celui où $S=T$, i.e. s'explicite comme une situation
% @ upper centre, diagram
\diagram{\begin{tikzcd}[ampersand replacement=\&,row sep=large]
 S \arrow[r,hook] \arrow[d,"p"'," {\text{surj}}"] \& \mathcal{L}\\
 I \arrow[d,"\varphi"'," {\text{surj.}}"] \& {}\\
 J \& {,}
\end{tikzcd}}
% slanted
% @ upper left, oblique note lines 1–8
\margin{Attention au cas où $\varphi$ est bijective, sans être un iso (deux structures d'ordre diff. sur un $=$ $I$, avec l'une plus fine que l'autre). C'est le cas d'un $\unsure{F|\mathcal F}\unsure{\leq|\preceq}G$, avec $F\ne G$, et \unsure{pourtant} $=$ relation \deleted{\unsure{d'éq}} $\unsure{\mathcal R|\mathscr R}_{F}=\mathcal R_{G}$ sur $|F|=|G|$.}
% @ beneath diagram
en posant $q=\varphi\circ p$.

% @ centre, heading
\uline{Proposition} \unsure{Soient} $F,G$ deux figures.

% @ centre, part a lines 1–3
\unsure{a)} Pour que $F$ raffine $G$, il \unsure{f} et s que ce soit une sous-figure d'une \deleted{\unsure{raff|suf}}\unsure{subdivision} $G'$ de \unsure{$F$|$G$}
% @ centre, relation
\[
 \unsure{:}\qquad F\leq G'\preceq G
\]

% @ lower middle, cancelled part b
\deleted{b) Pour que $F$ soit une subdivision de $G$, il f et s. que ce soit un raffinement, et qu'il soit maximal comme \unsure{Ainsi|Axiome}, \unsure{dès|donc} $F$}

% @ lower part, proof lines 1–5
\uline{Dém}. C'est \unsure{évid\(^{\mathrm{mt}}\)} suffisant, \unsure{Invers\(^{\mathrm{mt}}\)}, Supposons $F\ll G$, pour trouver $G'$, on prend $G'=G$ si $F$ est déjà une subd. i.e. $|F|=|G|$. \deleted{Sinon,} Dans le cas général, Soit
% @ bottom, two-line formula
\[
\begin{aligned}
 G'&=\bigl\{A\in\unsure{\mathfrak P|\mathscr P|\mathbb P}(L)\mid A\in F\text{ ou }A=\deleted{\unsure{g}}\{x\},\text{ où }x\in |F|\unsure{\setminus}\,\unsure{G|F}\unsure{,}\\
   &=\unsure{F}\coprod_{x\in |F|\setminus |G|}\{\{x\}\}
\end{aligned}
\]
\scanpage{27}\mspage{26}
% @ top, cancelled line
\deleted{\unsure{Alors|Ainsi} $\unsure{F|\mathcal F}$}

% @ top, paragraph lines 1–4
Ainsi, les relations d'ordre $\ll$ et $\preceq$ de, $F$ se déterminent mutuellement, en termes de $\leq$. Peut-on de $=$ exprimer $\leq$ en termes de $\ll$ et de $\preceq$~?
% @ top left, beside last line
\margin{??}

% @ upper middle, heading
\uline{NB} On a
% @ upper middle, implication
\[
 F\leq G\text{ et }F\preceq G\quad\Longrightarrow\quad F=G
\]
% @ middle, argument lines 1–6
\unsure{Il} faudra que ceci soit vrai dans \deleted{\uline{\unsure{toute}}} \unsure{contrée|contexte} -- via la condition que si $F\leq G$, $F\ne G$, $\exists\,H$ \inserted{\unsure{figure|ferm.}} \uline{non vide} $H$ disjointe de $F$, $H\ll G$ (et si on avait $F\preceq G$, $H$ serait disjointe de $G$ par C5, donc \unsure{OK|on} est \unsure{ramené} à ceci. \deleted{$H\ll G$,}
% @ middle right, implication with gloss
\[
 H\ll G,\quad
 \underset{\inserted{\substack{\text{(i.e. }H\text{ comp. à }G\text{ et}\\H\cap G=\varnothing}}}{H\text{ disjointe de }G}
 \quad\Longrightarrow\quad H=\varnothing_{F}
\]
% @ lower middle, proof lines 1–3
En effet, $H\leq G'\preceq G$, \deleted{\unsure{$i.e.$ $R_G$ $|H|$}} et $H$ disjointe de $G$ signifie $H$ disjointe de $G'$ \inserted{par \unsure{C5}}, donc $H=\varnothing$ OK~!)

% @ lower middle, cancelled heading
\deleted{\uline{Famille \unsure{intersection|intersections}}}

% @ lower part, question lines 1–2
Qu'en est-il des autres axiomes des \unsure{contrées|catories|catégories}~:

% @ foot, C5
\uline{C5} Si $F\preceq G$, $H$, alors $H$ dis. de $F$ ssi disj. de $G$ trivial, car $|F|=|G|$ et disjonction de $H$ avec $H'$ se voit sur les supports par $|H|\cap|H'|=\varnothing$

% @ foot, C6
\uline{C6} \deleted{\unsure{R|S}}\unsure{Subdivision} induite, on la verra \unsure{tantôt}

% @ foot, C7
\uline{C7} Ext. d'une subd. OK on le vérifiera \unsure{de tte façon|tantôt infra|tantôt itou}.

% @ bottom edge, short line
\uline{\unsure{C$\varnothing_0$|C$\varphi_0$}} trivial
\scanpage{28}\mspage{27}
% @ top, cancelled proposition lines 1–2
\deleted{\deleted{\uline{Proposition}} Soient $\unsure{F|\mathcal F},G$ deux figures \unsure{ens.} de $\mathcal L$. \deleted{\unsure{Cons}} Soit \deleted{$\unsure{K|H}$}}
% @ top, cancelled set
\[
 \deleted{K\deleted{\unsure{I|J}}=\{(A,B)\in F\times G\mid A^\circ\cap B^\circ\ne\varnothing\}}
\]
% @ upper part, cancelled connector
\deleted{de sorte que}
% @ upper part, cancelled coproduct
\[
 \deleted{|F\unsure{\cap\,\mid|\mid\cap\mid}G|=\coprod_{(A,B)\in \deleted{\unsure{I|J}}K} A^\circ\cap B^\circ}
\]
% @ middle, cancelled argument lines 1–5
\deleted{et que les \deleted{relations} $A^\circ\cap B^\circ$ $((A,B)\in K)$ forment les classes de la rel. d'équi. sur $|F|\cap|G|$ borne sup des relations \inserted{d'équi.} \deleted{\unsure{H}}induites par $\unsure{\mathcal R|\mathscr R}_F$ et $\mathcal R_G$. Considérons \deleted{la famille} \deleted{des} pour $k=(A,B)$}
% @ middle, last cancelled formula
\[
 \deleted{\deleted{\unsure{A_k}}\,C_k=A\cap B}
\]
% @ lower middle, first diagram
\diagram{\begin{tikzcd}[ampersand replacement=\&,column sep=small]
 \&\mathcal L\& {}\\
 S\arrow[ur,phantom,"\rotatebox{43}{$\subset$}" description]
  \arrow[d,"{\varphi\ \text{surj}}" ]\& {}\&
 \deleted{\unsure{T}}S'_{\deleted{\unsure{k|n}}}
  \arrow[ul,phantom,"\rotatebox{-42}{$\supset$}" description]
  \arrow[d,"{\varphi'\ \text{surj}}"]\\
 I\& {}\& I'\& {I,\ J\ \text{ens. ord.}}
\end{tikzcd}}

% @ lower part, introduction
Considérons
% @ lower part, second diagram
\diagram{\begin{tikzcd}[ampersand replacement=\&,row sep=small,column sep=small]
 {} \& {} \& S\arrow[dr,hook]\arrow[dddl,"\varphi"{swap,near start}] \& {}\\
 S\cap T\arrow[urr,hook]\arrow[drr,hook]\arrow[ddd,dash,"\unsure{\psi|Y}"'] \& {} \& {} \& L\\
 {} \& {} \& T\arrow[ur,hook]\arrow[dd,"\varphi'"] \& {}\\
 {} \& I \& {} \& {}\\
 \deleted{\unsure{J|I}}K\arrow[ur]\arrow[rr] \& {} \& I' \& {}
\end{tikzcd}}
% @ foot, image and set
avec
\[
 K=\operatorname{Im}\bigl(S\cap T\xrightarrow{(\varphi,\varphi')} I\times J\bigr)
 =\bigl\{(i,\unsure{j'|j})\in I\times I'\mid A_i^\circ\cap \deleted{\unsure{B}}\unsure{A_{i'}^{\prime\circ}|B_{j'}^{\prime\circ}}\ne\varnothing\bigr\}
\]
% @ bottom left, pair and boxed formula
si $j=(i,i')$, on aura $\boxed{A_j=A_i\cap A'_{i'}}$
% @ bottom right, two definitions
où on pose
\[
 \begin{aligned}
 \unsure{A_i'|A_i}&=\varphi^{-1}(I_{\leq i})\\
 \deleted{\unsure{A}}\unsure{B_j'|A_j'}&=\varphi^{\prime-1}(I'_{\leq j})
 \end{aligned}
\]
\scanpage{29}\mspage{28}
% @ top, false start
\deleted{Alors}

Ainsi

% @ top, proposition lines 1–3
\uline{Proposition} On a construit un Inf de $\unsure{F|\mathcal F}$ et $G$ pour $\ll$, on \deleted{\unsure{y|l}}l'appelle la ``figure intersection'' $F.G$.

% @ upper middle, warning lines 1–2
Attention, ce \uline{n'est pas} \inserted{\unsure{en général}} la figure $F\cap G$, \deleted{\unsure{intersection|inter.}} en se rappelant $F,\deleted{\unsure{F}}G\subset\unsure{\mathfrak P|\mathscr P|\mathbb P}(L)$, \deleted{$G\subset\mathfrak P(L)$}

% @ middle, paragraph lines 1–3
Nous n'allons considérer cette $F\cap G$ que dans le cas où $F$ et $G$ compatibles, mais alors les deux coïncident~!

% @ middle, paragraph lines 1–4
Donc la bonne notion d'intersection est le $F.G$. En fait, ça marche pour une famille quelconque de figures, et
% @ middle lower, support formula
\[
 \left|\bigwedge_{i\in I} F_i\right|=\bigcap_{i\in I}|F_i|\quad .
\]

% @ lower part, paragraph lines 1–4
Cette propriété d'existence d'un Inf pour la relation $\ll$ \unsure{s'apparente} \deleted{\unsure{la propriété|la présentation}} à une propriété de ``\uline{\unsure{modération|médiation}}''.

% @ foot, cancelled statement lines 1–3
\deleted{\uline{Propriété} Si $F'\preceq F$ \deleted{\unsure{, et des}}, alors $F'.G\preceq F.G$ (\unsure{Immédiat|Immédit|immédite} sur supports)~: Subdivision \uline{induite} sur un raffinement.}
\scanpage{30}\mspage{29}
% @ top, proposition heading
\uline{Proposition} Soient ${\unsure{F|\mathcal F}}',F,G$ trois figures.
% @ top, three implications
\[
\begin{aligned}
 \text{a)}\quad F'\ll F&\Longrightarrow F'.G\ll F.G\\
 \text{b)}\quad F'\preceq F&\Longrightarrow F'.G\preceq F.G\\
 \text{c)}\quad F'\leq F&\Longrightarrow F'.G\leq F.G\quad *
\end{aligned}
\]
% @ upper right, beside a
\margin{(trivial par \unsure{déf.} de $\ll$ comme Inf)}
% @ upper right, beside b
\margin{(\unsure{résulte} de a), en regardant supports)}
% @ upper left, starred gloss
\margin{$*$ \uline{\unsure{explicitation|expliciter|explicitons}}}
% @ upper part, starred formula
\[
 F'.G=\{A\in F.G\mid A\subset |F'|\}
\]

% @ upper part, proof lines 1–2
Prouvons c. Il faut prouver que \deleted{$F'.G$} $F'.G\subset F.G$, et est fermé dans $F.G$.

% @ middle, proof lines 1–5
\deleted{\unsure{or}} Or $F'.G$ est formé des $A\cap B$ $(A\in F',B\in G)$ tels que $A^\circ\cap B^\circ\ne\varnothing$, et $F.G$ des $A\cap B$ $(A\in F,B\in \unsure{G|L})$ tels que $A^\circ\cap B^\circ\ne\varnothing$, et dans les deux cas le \unsure{sens} de $A^\circ$ (\unsure{p.r.} à $F$, \unsure{dans|ds} $F'$) est le $=$, donc $F'.G\subset F.G$. Prouvons que c'est fermé i.e.
% @ centre, first membership and annotations
Si \deleted{$(A\cap B\qquad A,A'\in F',$}
\[
 \underset{\substack{\rotatebox{-90}{$\in$}\\ F'}}{A}\cap
 \underset{\substack{\rotatebox{-90}{$\in$}\\ G}}{B}\in F'.G\qquad
 \text{i.e. } A^\circ\cap B^\circ\ne\varnothing
\]
% @ centre, second membership and annotations
\[
 \underset{\substack{\rotatebox{-90}{$\in$}\\ F}}{A'}\cap
 \underset{\substack{\rotatebox{-90}{$\in$}\\ G}}{B'}\in\underset{\text{\inserted{i.e. $A'^\circ\cap B'^\circ\ne\varnothing$}}}{F.G},\qquad
 \text{et } A'\cap B'\subset A\cap B\ \deleted{\unsure{\text{i.e.}}},\quad\text{alors}
\]
% @ middle lower, conclusion lines 1–4
$A'\cap B'\in F'.G$ i.e. $A'\in F'$. Or on sait que (d'après la théorie de $F.G$) que $A'\cap B'\subset A\cap B\Longrightarrow A'\subset A$, $B'\subset B$, d'où $A'\in F'$ puisque $F'$ fermé dans $F$.

% @ lower left, oblique explanation
% slanted
\margin{\uline{Expliciter}~: soit
\diagram{\begin{tikzcd}[ampersand replacement=\&,column sep=small,row sep=small]
 G\arrow[r,phantom,"\leq" description]\&F\arrow[d,phantom,"\rotatebox{90}{$\preceq$}" description]\\
 {}\&F'
\end{tikzcd}}, d'où $G'$ avec $G'\preceq G$, $G'\leq F'$. Alors on a $G'=\{A'\in F'\mid A'\subset |G|\}=\varphi^{-1}(G)$, où $\varphi:F'\to F\ldots$}
% @ lower right, note on uniqueness
\margin{\unsure{Un peu vif}, on trouve par $F'.G=G'$ un $G'$ tel que
\diagram{\begin{tikzcd}[ampersand replacement=\&,column sep=small,row sep=small]
 F'\arrow[r,phantom,"\preceq" description] \arrow[d,phantom,"\rotatebox{90}{$\leq$}" description]\& F\arrow[d,phantom,"\rotatebox{90}{$\leq$}" description]\\
 G'\arrow[r,phantom,"\preceq" description]\& G
\end{tikzcd}}
mais il faut vérifier encore l'\uline{unicité}}
% @ lower right, second note
\margin{Cela amène à dégager a) $P\ll Q\ll R$ et $P\preceq R\Rightarrow P\preceq Q,\ Q\preceq R$ b) $P\ll Q\ll R$ \unsure{et} $P\leq R$ \deleted{$\Rightarrow P\unsure{\leq|<}$} $Q\leq R$ $\Rightarrow P\leq Q$}
% @ lower part, corollary
\uline{Cor}. L'axiome C6 est \unsure{valable|valide}.

% @ lower part, introduction to C7
Vérifions quand même C7.
% @ lower part, order diagram and existence
\diagram{\begin{tikzcd}[ampersand replacement=\&,column sep=small,row sep=small]
 G\arrow[r,phantom,"\leq" description]\arrow[d,phantom,"\rotatebox{90}{$\preceq$}" description]\&F\& {\exists\ \text{subdivision }F'\text{ de }F\text{ qui}}\\
 G'\& {}\&{\text{induise }G'.}
\end{tikzcd}}
% @ bottom, construction
On prend $F'=G'\cup\deleted{\unsure{(F.G)}}(F\setminus G)$ (réunion disjointe)
% @ bottom, verification items
Il faut vérifier a) c'est une figure b) \unsure{qui|que} \deleted{\unsure{raffine}}\inserted{\unsure{subdivise}} $F$ et c) qui induit $G'$ sur $G$.
\scanpage{31}\mspage{30}
% @ top, part a
a) ${\unsure{F|\mathcal F}}'$ est une figure. \deleted{i.e.} Notons que \deleted{$|F'|=$}
% @ top, support with braces
\[
 |F'|=\underbrace{|G'|}_{|G|}\cup\underbrace{|F\setminus G|}_{\supset |F|\setminus |G|}\deleted{\unsure{=|\supset}}\supset |F|
\]
% @ upper part, connector
donc, comme $|F'|\subset |F|$, on a
% @ upper part, support equality
\[
 |F'|=|F|
\]
% @ upper part, argument
Il faut prouver que $\forall x\in |F'|=|F|$, $\exists$ plus petit \unsure{élémt|élément|élt} de $F'$ qui le contienne. Distinguons deux cas

% @ upper middle, first case
1$^\circ$) $x\in |G|=|G'|$, soit $A\in G'\deleted{\unsure{\circ}}$ tel que $x\in A^{\circ(G')\deleted{\unsure{)|?}}}$. Si $B\in F'$ et \deleted{\unsure{B}} $x\in B$, prouvons $A\subset B$. Si $B\in G'$, c'est évident. Si $B\in F\setminus G$, alors \deleted{$B$ contient} $A^{\circ(G')}\cap B\ne\varnothing$ donc $A^{\unsure{\bullet|\circ}}\subset B$ OK

% @ middle, second case
2$^\circ$) $x\in |\deleted{\unsure{G}}F|\setminus |G|$, alors les $\deleted{\unsure{B}}A\in F'$ qui contiennent $x$ ne sont pas dans $G'$, ils sont donc dans $F\setminus G$. D'autre part, l'unique \unsure{élémt|élément|élt} $A$ de $F$ tel que $x\in A^{\circ(F)}$ n'est pas dans $G$ puisque $x\notin |G|$, il est donc dans $F\setminus G$, et \deleted{\unsure{si $x\in A$}} c'est le plus petit él. de $F$ (et à fortiori de $F\setminus G$) qui contient $x$, OK~!

% @ lower middle, part b
b) $F'$ \deleted{\unsure{raffine}}\inserted{subdivise} $F$. Comme $|F'|=|F|$, il faut voir \unsure{seult} que $F'$ raffine $F$. Donc \deleted{que \unsure{tout}} que
% @ lower part, cancelled attempt
\deleted{1$^\circ$) \deleted{Soit} $A\in F'$, \unsure{et} l'ens $F^A$ des $B\in F$ tels que $A\subset B$ \unsure{est contenu dans} un plus petit \inserted{\unsure{a un plus petit}} \unsure{élémt|élément|élt} de $F$. C'est évident si $A\in F$ c'est évident. Si}
\deleted{\unsure{str.} de de $F'$ raffine $F$. Si}

% @ lower part, restated first condition
1$^\circ$) \deleted{Soit} $A\in F'$, l'ens $F^A=\{B\in F\mid \deleted{\unsure{B}}A\subset B\}$ a un plus petit \unsure{élémt|élément|élt}. C'est évident si $A\in F$ \deleted{$(B=A)$} (on prend $B=A$) et si $A\in G'$, c'est vrai puisque $G'\ll F$.

% @ bottom, second condition
2$^\circ$) Si $A\in F'$, $B\in F$, alors $A\cap B$ est \deleted{\unsure{$\mathcal R_{F'}$-satu.|$\mathscr R_{F'}$-satu.}}\inserted{fermé dans $F'$}, i.e. réunion de strates de $F'$. Si $A\in G'$, comme $G'\ll F$, $A\cap B$ est réunion de strates de $G'$, qui sont aussi des strates de $F'$.
\scanpage{32}\mspage{31}
% @ top, continuation lines 1–3
Si $A\in \unsure{F|\mathcal F}\setminus G$, $A\cap B$ est réunion de strates de $F$. \deleted{\unsure{Reste}} reste à voir que une strate \inserted{\unsure{$G$|$\in$}} de $F$ est réunion de strates de $F'$. Soit $A$ cette strate de $F$. \deleted{\unsure{On}}
% @ top, cancelled union
\[
 \deleted{A=\bigcup_{\substack{B\in F\\B\leq F}}B^\circ}
\]
% @ upper part, two cases
Si $A\notin G'$, alors $A\in F\setminus G'\subset F'$, et on gagne. Si $A\in G$, alors (comme $G'\preceq G$) $A$ est réunion de strates de $G'$, qui sont des strates de $F'$, et on gagne encore.

% @ left beside corollary
% slanted
\margin{\unsure{Ultra} faux}
% @ upper middle, corollary heading
\uline{Corollaire} \unsure{Le|La} \deleted{raffinement}\inserted{subdivision} précédente $F'$ de $F$ est caractérisée par \inserted{chacune} des propriétés équivalentes suivantes~:

% @ centre, crossed condition a
\deleted{a) Pour toute subdivision $F''$ de $F$ qui induit $G'$ sur $G$, on a $F''\preceq F'$ -- i.e. $F'$ est \deleted{\uline{maximal}}\inserted{un plus grand \unsure{élémt|élément|élt}} pour $\preceq^{*}$, dans l'ens des subdivisions de $F$ qui induisent $G'$ sur $G$.}
% @ left of a, starred gloss
% slanted
\margin{$*$ ou pour $\ll$, cela revient au $=$ \unsure{parmi|pour} des subdivisions d'une $=$ $F$}
% @ centre left, long oblique note
% slanted
\margin{les conditions \deleted{\unsure{a)}} b) c) d) sont équivalentes, \unsure{et} \uline{\unsure{impliquées}} par a), \unsure{lesquelles|laquelle} \unsure{caractérisent|caractérise} la subdivision construite \unsure{précédemment}. Mais ces trois \unsure{cond.} (b,c,d) sont \unsure{loin} d'impliquer a)}
% @ lower middle, condition b
b) L'application \unsure{can.} associée à $F'\preceq F$
\[
 \varphi:F'\longrightarrow F
\]
induit une bijection
% @ lower middle, bijection
\[
 F'\setminus G'\xrightarrow{\sim}F\setminus G.
\]
% @ right, NB
\margin{NB. c'est \unsure{tjrs|tojrs} surj. donc c'est l'\uline{injectivité} qu'il faut prouver ici}
% @ lower part, condition c
c) \deleted{Les strates ouvertes de $F'$ qui ne sont} $\forall A\in F\setminus G$, $A^\circ$ est une strate ouverte de $F'$.

% @ lower part, condition d
d) $F'\supset F\setminus G$, i.e. $\forall A\in F\setminus G$, on a $A\in F'$.

% @ foot, proof heading
\uline{Dém} \deleted{\unsure{On}} \unsure{On} va prouver \unsure{p. ex}
% @ foot, cycle of implications
\[
 \text{d)}\Longrightarrow\text{b)}\Longrightarrow\text{c)}\Longrightarrow\text{a)}\Longrightarrow\text{d)}
\]
% @ bottom, first implication
\deleted{\uline{\unsure{Dém}}} $\text{d)}\Rightarrow\text{b)}$ Si $(F\setminus G)\subset F'$, alors $\varphi\mid F\setminus G$ est l'\deleted{\unsure{identité}}\unsure{inclusion} \deleted{(trivial)} \inserted{\unsure{can.} $F\setminus G\hookrightarrow F$}, \deleted{donc} donc il reste à prouver
% @ bottom edge, equality
\[
 F\setminus G=F'\setminus G'.
\]
\scanpage{33}\mspage{32.}
% @ top, 1°)
\deleted{\unsure{R|P}} 1°) $\unsure{F|\mathcal F}\setminus G\subset F'\setminus G'$ i.e. Tout $A\in F\setminus G$ est $\notin G'$. En effet, si on avait $A\in G'$, on aurait $|A|\subset|G'|=|G|$ donc \deleted{$A\in$} (comme $G\leq F$) $A\in G$.

% @ left edge, along 1°) and 2°)
% margin note written sideways
\margin{\unsure{voir} \unsure{une|sa} \unsure{simplification}, cf plus bas}

% @ upper middle, 2°)
2°) $F'\setminus G'\subset F\setminus G$ Soit $A'\in F'\setminus G'$, prouvons $A'\in F\setminus G$, et d'abord $A'\in F$. \deleted{On a} Soit $A$ le plus petit él. de $F^{A'}$ donc $\underline{\rule[-.15ex]{.35pt}{.55ex}\,A'\subset A^{\deleted{\unsure{\prime|\ast|?}}},\ A'^{\circ(F')}\subset A^{\deleted{\unsure{\prime|\ast|?}}\circ(F)}\,\rule[-.15ex]{.35pt}{.55ex}}$ prouvons que $A'=A$ (d'où $A\in F$). \deleted{On a}

% @ middle, Je dis
Je dis que $A\in F\setminus G$ i.e. $A\notin G$, sinon on aurait $A\subset|G|=|G'|$ donc $A'\in G'$ puisque $A'\in F'$ et $G'\leq F'$, absurde. \deleted{Mais} Donc $A\in F\setminus G$, donc $A\in F'$, \deleted{\unsure{et}} Il en résulte que $A^{\circ(F)}=A^{\circ(F')}$ (car $F'\preceq F$\inserted{il suffirait $F'\ll F$}\deleted{)} cf lemme plus bas) et \deleted{\unsure{bien}} \deleted{$\unsure{F|\in}$} $A,A'\in F'$ tels que $A'^{\circ(F')}\subset A^{\circ(F')}$, ce qui prouve $A=A'$.

% @ lower middle, Donc
Donc \unsure{on a} prouvé que $\deleted{\unsure{F|A}}A'\in F'\setminus G'\Rightarrow A'\in F$. Prouvons qu'on a $A'\notin G$ (d'où $A'\in F\setminus G$\deleted{)}, qed). Sinon on aurait $A'\in G$, on aurait $A'\subset |G|=|G'|$, d'où (comme $A'\in F'$ et $G'\leq F'$) $A'\in G'$, absurde.

% @ lower third, cancelled b) to c)
\deleted{b) $\Rightarrow$ c) \inserted{On suppose b) i.e. $\varphi$ induit $F'\setminus G'\overset{\sim}{\longrightarrow}F\setminus G$} Soit $A\in F\setminus G$, prouvons que $A^{\circ(F)}$ est une \uline{strate ouverte \deleted{dans \unsure{la}} de $F'$}. Soit $A'\in F'\setminus G'$ l'unique élément \deleted{de} tel que $\varphi(A')=A$. Prouvons $A'^{\circ(F')}=A^{\circ(F)}$.}

% @ bottom, b to c
b $\Longleftrightarrow$ c \deleted{Posons} On a $|F|\setminus|G|=|F'|\setminus|G'|$, soit $U$, Posons $F\setminus G=I^{\unsure{\ast|k|+}},\ F'\deleted{\unsure{\times|\setminus}}\unsure{\setminus}G'=I'^{\ast}$ (parties ouvertes des ens d'indices ordonnés $I$, \deleted{$\unsure{A|\Delta}$} $(=F)$, $I'(=F')$. On a
